UX evaluation is costly, slow, and often deferred until a product has enough users to justify studies.
For early-stage teams, this creates a gap: product decisions must be made under uncertainty,
while regressions and UX friction can accumulate silently until they become expensive to unwind.
In parallel, large language model (LLM) tooling has made it feasible to generate structured, role-conditioned feedback at scale,
but the scientific and practical question remains: \emph{what does such synthetic feedback mean}, and how can
we validate it without over-claiming?

This paper studies a pragmatic middle ground: we use LLM-driven \emph{micro-simulations} to generate
structured UX feedback from versioned prompts, and validate the resulting friction signals against
multiple public \emph{proxy datasets} (app reviews, customer support tweets, open-source software issues).
Rather than treating proxies as ground truth for task success, we treat them as noisy distributions of
friction themes that can be compared to simulation outputs under controlled configurations.

\paragraph{Research questions and hypotheses.}
We focus on three practical questions:
(RQ1) do simulation outputs align with friction distributions observed in public proxy corpora?
(RQ2) which proxy-alignment method is most robust (lexical vs multilingual embeddings, and which overlap metric)?
and (RQ3) are conclusions stable under sampling variability and prompt/taxonomy changes?
Our main hypotheses are that embedding-based overlap will outperform lexical baselines under weighted-Jaccard (H1),
that top-$k$ Jaccard saturates and can overstate alignment when $k$ grows (H2),
and that aggregate alignment estimates are sufficiently stable to support iterative decision-making when reported with bootstrap intervals (H3).

\paragraph{Contributions.}
We make three contributions:
(1) a reproducible UX micro-simulation pipeline with structured outputs and versioned inputs;
(2) a proxy-validation protocol with alignment metrics (top-$k$ and weighted Jaccard),
including lexical and multilingual embedding baselines; and
(3) an artifact-first evaluation suite with bootstrap confidence intervals and error analyses, designed
to be upgraded with stronger paid models only at the end of the iteration loop.

\paragraph{Value to practitioners.}
For product and engineering teams, the output of our pipeline is not a single ``UX score,'' but a set of actionable signals:
(i) a ranked set of friction themes to prioritize in the backlog, (ii) a compact summary of where the current experience
is likely to fail (by locale/domain proxy), and (iii) an auditable, versioned trail of prompts, taxonomies, and artifacts
that makes changes explainable and reviewable.
In other words, we turn otherwise unstructured, costly-to-collect ``user feedback'' into a reproducible iteration loop
that can be consumed by stakeholders who need to decide \emph{what to fix next} and \emph{what risks are acceptable} in early stages.
